\documentclass[11pt]{article}
\usepackage[utf8]{inputenc}
\usepackage{geometry}
\geometry{a4paper,margin=1in}
\usepackage{hyperref}
\usepackage{booktabs}
\usepackage{longtable}
\title{RMT-LLM Laboratory — Experiment Report}
\author{Iskhak Hamzatovich Isaev}
\date{\today}
\begin{document}
\maketitle
\section{Detailed Results with Explanations}
 Experiment Metadata

- **Name:** `Spectral Signature`
- **Language:** `unknown`
- **Version:** `unknown`
- **Scenario:** `default`

__

 Key Metrics

 |  Metric  |  Value  |
 | --- | --- |
 |  n_layers_analyzed  |  4  |
 |  max_lambda_max  |  0.012519393000410352  |
 |  min_lambda_min  |  0.00017214117106267862  |
 |  mean_spectral_gap  |  0.005858584940683028  |

 RMT Spectral Analysis

 |  Quantity  |  Value  |  Interpretation  |
 | --- | --- | --- |

 Reasoning Trace (Hidden Chain-of-Thought)

 |  Metric  |  Value  |
 | --- | --- |

 Interpretation

The **Marchenko-Pastur (MP) law** describes the bulk distribution of eigenvalues of large random covariance matrices. When the largest empirical eigenvalue exceeds the MP upper bound (`mp_upper`), this signals structured (non-random) information — i.e. the model has 'detected' a fact. Conversely, when eigenvalues stay inside the MP bulk, the model is generating creative or hallucinated content.

The **N_crit threshold** is the autoregressive token count beyond which spectral collapse makes hallucination mathematically inevitable. Below N_crit the model can still self-correct; above it, the Caputo fractional dynamics dominate and the model enters the **'utility trap'** described in the project's main monograph. The observed behavior verifies (or falsifies) the RMT-LLM framework's prediction.

The **reasoning trace** exposes the model's hidden chain-of-thought. A `mean_deception` > `mean_honesty` indicates the model is internally planning to mislead the user, while `filter_bypass_count > 0` confirms that protective filters fire only at output time, not at reasoning time — exactly as reported in the source news article.
\section{Full Task Launch Logs}
\begin{verbatim}
\end{verbatim}
\end{document}
